Large Language Models (LLMs) is indispensable in today's applications, but their inference procedure—generating responses by breaking text into smaller pieces and processing them using a memory-heavy element named Key-Value (KV) cache—requires a lot of computational resources, especially when memory is limited. This paper treats LLM inference optimization as a multi-stage online scheduling problem, where prompts arrive sequentially and the incremental expansion of the KV cache during inference renders conventional scheduling algorithms ineffective. To address this challenge, we develop a fluid dynamics approximation to establish a tractable benchmark, providing insights for devising effective scheduling algorithms. Building upon this foundation, we introduce the Nested Waiting for Accumulated Inference Threshold (Nested WAIT) algorithm. This algorithm constructs a hierarchical framework comprising multiple segments, each defined by distinct thresholds, to effectively manage the random prompt arrivals with unknown output lengths. Theoretical analysis shows both algorithms have near-optimal performance compared with the fluid benchmark under heavy traffic limit, balancing throughput, latency, and Time to First Token (TTFT). Numerical experiments conducted with the Llama-7B model on an A100 GPU, utilizing real-world datasets, demonstrate that our approach achieves superior throughput and reduced latency compared to widely adopted baseline methods such as vLLM and Sarathi. This research bridges operations research and machine learning, presenting a theoretically grounded framework for the efficient deployment of large language models under memory constraints.